\documentclass[12p, paper=a4, leqno, colorinlistoftodos]{article}
\usepackage{wrapfig}
\usepackage[utf8]{inputenc}
\usepackage[ngerman]{babel}
\usepackage{amsmath}
\usepackage{amsfonts} 
\usepackage{amssymb}
\usepackage{enumerate}
\usepackage{graphicx}
\usepackage{wrapfig}

\usepackage[parfill]{parskip}
\usepackage{subcaption}
\usepackage{fix-cm}
\usepackage{transparent}
\usepackage{color}
\usepackage{xcolor}
\usepackage{tocloft, url}
\usepackage{hyperref}
\usepackage{etoolbox}
\usepackage{fancyhdr}
\usepackage[german]{fancyref}
\usepackage{lipsum}	
\usepackage{sectsty}
\usepackage{pdfpages}
\usepackage{tcolorbox}
\usepackage{todonotes}


\usepackage{enumitem}
\usepackage[bottom]{footmisc}

\usepackage{forest}
\usepackage{float}
\usepackage{longtable}
\usepackage[normalem]{ulem}
\usepackage{multirow}
\usepackage[font=small,labelfont=bf,tableposition=top]{caption}
\usepackage{tikz}
\usepackage{calc}
\usepackage{geometry}
\usepackage{listings,lstautogobble}
\usepackage{algorithm}
\usepackage{algorithmicx}
\usepackage[noend]{algpseudocode}

\usepackage{array}
\usepackage{booktabs}
\usepackage{xcolor,colortbl}
\usepackage{setspace}
\usepackage{titlesec}
\usepackage[nottoc,notlot,notlof,numbib]{tocbibind}

\usepackage{cite}


\hypersetup{
	colorlinks,
	linkcolor={black},
	citecolor={blue!50!black},
	urlcolor={blue!80!black}	
}

% BEGIN COLORS
\definecolor{CBlue}{RGB}{1,0,119}
\definecolor{CBlack}{rgb}{0.0,0.0,0.0}
\definecolor{CLightGrey}{rgb}{0.664,0.664,0.664}
\definecolor{LightOrange}{rgb}{1.0,0.9,0.63}
\definecolor{DarkOrange}{rgb}{1.0,0.8,0.22}
\definecolor{LightCyan}{rgb}{1.0,0.9,0.63}
\definecolor{CWhite}{RGB}{0,0,0}
\definecolor{CLink}{rgb}{0.5,0.0,0.0}

\newcommand{\CTitle}{CBlue}
% BEGIN Table row/column colors{102,217,239}
\definecolor{CEven0}{RGB}{208,238,244}
\definecolor{CUneven0}{RGB}{255,255,255}
\definecolor{CEven1}{rgb}{1.0,1.0,1.0}
\definecolor{CUneven1}{rgb}{1.0,1.0,1.0}
% END Table row/column colors
% END COLORS


% BEGIN NEWENVIRONMENTS
\newenvironment{packed_itemize}
{\begin{itemize}
		\setlength{\itemsep}{0pt}
		\setlength{\parskip}{0pt}
		\setlength{\parsep}{0pt}
	}{\end{itemize}}

\newenvironment{packed_description}
{\begin{description}
		\setlength{\itemsep}{0pt}
		\setlength{\parskip}{0pt}
		\setlength{\parsep}{0pt}}
	{\end{description}}

\newenvironment{packed_enumerate}
{\begin{enumerate}
		\setlength{\itemsep}{0pt}
		\setlength{\parskip}{0pt}
		\setlength{\parsep}{0pt}
	}{\end{enumerate}}

% END NEWENVIRONMENTS


% BEGIN NEWCOMMANDS
% tables
\newcommand{\tabularPaddingSmall}{\renewcommand{\arraystretch}{0.75}}
\newcommand{\tabularPaddingDefault}{\renewcommand{\arraystretch}{1}}
\newcommand{\tabularPaddinglarge}{\renewcommand{\arraystretch}{1.25}}
\newcommand{\tabularPaddingLarge}{\renewcommand{\arraystretch}{1.5}}
\newcommand{\tabularPaddingLARGE}{\renewcommand{\arraystretch}{2}}
\newcommand{\tabularPaddingBot}{\vspace{0.5cm}}

\newcommand{\tabitem}{~~\llap{\textbullet}~~}
\newcommand{\coloredHline}[1]{\arrayrulecolor{#1}\hline}
\newcommand{\minitab}[2][l]{\begin{tabular}{#1}#2\end{tabular}}

% quoting
\newcommand{\quoteSource}[2]{\glqq{#1}\grqq, \textit{#2}}
\newcommand{\quoteQuick}[1]{\glqq{#1}\grqq}
\newcommand{\quoteAuthor}[2]{\quoteQuick{#1}, von \textit{#2}}
\newcommand{\quoteCite}[2]{\quoteQuick{\textit{#1}}#2}
\newcommand{\srcSelf}{\\\text{Quelle: Eigene Darstellung}}
\newcommand{\citePage}[2]{\cite{#1} Seite: {#2}}
\newcommand{\citeEquation}[1]{\text{\hspace{1cm}{#1}}}

% paragraphs
\newcommand{\parTiny}{\vspace{2mm}\\}
\newcommand{\parSmall}{\\~\par}
\newcommand{\parNoIndent}{\\\noindent}
\newcommand{\parWrapfig}{\\\text{ }\\}

% math


% 
\newcommand{\emphFold}[1]{\textbf{#1}}

% page layout
\newcommand{\geometryNormal}{\newgeometry{top = 2cm, left=2.5cm,bottom=2cm, right = 4cm}}
\newcommand{\geometryToc}{\newgeometry{top = 3cm, left=3cm,bottom=3cm, right = 3cm}}

% todo notes
\newcommand{\todoLowp}[1]{\todo[color=green!40]{#1}}
\newcommand{\todoMedp}[1]{\todo[color=orange!40]{#1}}
\newcommand{\todoHighp}[1]{\todo[color=red!40]{#1}}
\newcommand{\todoFig}[1]{\missingfigure{#1}}

% titles
\colorlet{ctcolorchapternum}{black}

\newcommand\mysectionformatCore[1]{%
	\vspace{-3em}\parbox[b]{\dimexpr\textwidth-1em\relax}{\raggedright#1}}
\newcommand{\titleformatCore}
{
	\titleformat
	{\section}[display]%
	{\color{black}\Large\bfseries}%
	{\vspace{-18em}\raggedleft{%
			{\color{ctcolorchapternum}\fontsize{60}{60}{\selectfont\thesection}}%
		}\hspace{1em}%
	}%
	{0pt}%
	{\mysectionformatCore}%
	[{\vspace{-0.5em}%
		\rule{0.75\textwidth}{1pt}%
		\rule[-3pt]{0.25\textwidth}{4pt}}]
	
	\titleformat
	{\subsection}[block]%
	{\color{black}\large\bfseries}%
	{\thesubsection\text{ }}%
	{0pt}{}[]%
	
	\titleformat
	{\subsubsection}[block]%
	{\color{black}\normalsize\bfseries}%
	{\thesubsubsection\text{ }}%
	{0pt}{}[]%
	
	\titleformat{\paragraph}
	{\normalfont\normalsize\bfseries}{\theparagraph}{1em}{}
	\titlespacing*{\paragraph}
	{0pt}{3.25ex plus 1ex minus .2ex}{1.5ex plus .2ex}
}

\newcommand\mysectionformatAppendix[1]{#1}
\newcommand{\titleformatAppendix}
{
	\titleformat
	{\section}[display]%
	{\color{black}\Large\bfseries}%
	{}{0pt}{\mysectionformatAppendix}[]
}

% END NEWCOMMANDS

% BEGIN RENEW COMMANDS
% END RENEW COMMANDS

\newlistof{links}{lks}{List of Links}
\newcommand\externallink[1]{%
	\refstepcounter{links}%
	\footnote{\url{#1}}%
	\addcontentsline{lks}{links}{%
		\protect\numberline{\thelinks}%
		\protect\url{#1}}%
}
\newcommand\externallinkIL[1]{%
	\refstepcounter{links}%
	\addcontentsline{lks}{links}{%
		\protect\numberline{\thelinks}%
		\protect\url{#1}}%
}

\newcolumntype{L}[1]{>{\raggedleft\let\newline\\\arraybackslash\hspace{0pt}}m{#1}}
\newcolumntype{R}[1]{>{\raggedright\let\newline\\\arraybackslash\hspace{0pt}}m{#1}}
%\newcolumntype{C}[1]{>{\raggedcenter\let\newline\\\arraybackslash\hspace{0pt}}m{#1}}

\newcolumntype{C}[1]{>{\raggedcenter\hspace{0pt}}p{#1}}


\author{Matthias Elmar Gensheimer}
\title{EyeCandy3D Manual}
\date{\today}


% BEGIN HEADER
\newcommand{\headerNormal}
{
	\pagestyle{fancy}
	\fancyhf{}
	\lhead{\rightmark}
	\rhead{\thepage}
}

\newcommand{\headerEmpty}
{
	\fancyhf{}
	\renewcommand{\headrulewidth}{0pt}
	\lhead{}
	\rhead{}
}

\newcommand{\headerErklaerung}
{
	\pagestyle{fancy}
	\fancyhf{}
	\rhead{\thepage}
}
% END HEADER

% BEGIN FOOTER
\usepackage{scrlayer}
\DeclareNewLayer[
foreground,
foot,
hoffset=0pt,
width=\paperwidth,
contents={\parbox{\layerwidth}{\centering}}
]{PageMarkCentredToPage}
\RedeclarePageStyleByLayers{plain}{PageMarkCentredToPage}
% END FOOTER

% BEGIN TOC SETUP
\makeatletter
\patchcmd{\l@section}
{\hfil}
{\leaders\hbox{\normalfont$\m@th\mkern \@dotsep mu\hbox{.}\mkern \@dotsep mu$}\hfill}
{}{}
\makeatother
% END TOC SETUP

% BEGIN PAGE SETUP
\geometryNormal
\setlength{\extrarowheight}{3pt}
% END PAGE SETUP


\titleformatCore
\headerNormal

\begin{document}
	\pagenumbering{Roman} 
	
	\lstset{language=C++,
		keywordstyle=\color{blue},
		basicstyle=\ttfamily,
		commentstyle=\ttfamily\itshape\color{gray},
		stringstyle=\ttfamily,
		showstringspaces=false,
		breaklines=true,
		frameround=ffff,
		rulecolor=\color{black},
		frame=single,
		autogobble=true
	}
	
	\begin{titlepage}
		\begin{flushright}
			\includegraphics[width=0.5\textwidth]{resources/hs_kempten_logo}
		\end{flushright}
		\begin{center}
			\topskip0pt
			\vspace*{\fill}
					\LARGE\textbf{EyeCandy3D User Manual}\\		
		
			\LARGE{A 3D Scenegraph built on OpenGL}\\
			
			
			\vfill
			\Large{\textcopyright 2018 Matthias Gensheimer}\\
			\Large{All Rights Reserved.}
		\end{center}
		
	\end{titlepage}
	
	% Add an empty page
	\mbox{}
	\thispagestyle{empty}
	\newpage
	
	\pagebreak
	
	\geometryNormal
	\tableofcontents
	\geometryNormal
	\setstretch{1.30}
	
	\pagenumbering{arabic} 
	\section{Introduction}
	EyeCandy3D is a scene graph based on OpenGL. Every program consists of one \textbf{Application}, which holds a number of windows.
	Each window can hold a variable number of scenes, which can run simultaneously.
	
	
	
	\pagebreak
	\section{First steps}
		In this section I will guide you through the most important aspects of using this library.
		
		\subsection{Creating an application}
		An \textbf{Application} is the main part of this library.
		\begin{lstlisting}
			Application app;
		\end{lstlisting}
		It allows the creation of new windows, which will automatically visualized. The following function will create a new window and add it to the application:
		\begin{lstlisting}
			app.createWindow<WindowType>(...);
		\end{lstlisting}
		\textbf{WindowType} has to be an own \textbf{Window} implementation.
		
		\subsection{Creating a window}
		Custom windows must be derived from \textbf{ec::Window}. This window has to be instanced with the window creation function provided by an \textbf{Application}. 	
		\begin{lstlisting}
			class ExampleWindow : public ec::Window
			{
			public:
				explicit ExampleWindow(unsigned int windowWidth,
					unsigned int windowHeight,
					const std::string& title)
				// Own window implementation ...
			};
		
			ec::Application app;
			app.createWindow<ExampleWindow>(width, height, "Example Window", "id");
		\end{lstlisting}
		
		The dervied window  \textbf{must} support a constructor with the following signature:\\
		\textbf{MyWindow(unsigned int, unsigned int, const std::string\&)}
		
		
		
		\subsection{Creating a scene}
		Custom scenes must be derived from \textbf{ec::Scene}. By registering it with the scene manager in the associated window, it can receive updates.	
		
		\begin{lstlisting}
			class ExampleScene : public ec::Scene
			{
			public:
				explicit ExampleScene(const std::string& name, 
					ec::Window* window);
			
				// Own window implementation ...
			};
			
			// In ExampleWindow:
			void ExampleWindow::initScenes()
			{
				// Create new window
				auto exampleScene = new ExampleScene("example", this);
				
				// Register the new window, so it receives updates
				m_sceneSystem.registerScene(exampleScene);
			
				// ...
			}
		\end{lstlisting}
		
		\subsection{Resource registry}
		Built in types:
		\begin{itemize}
			\item 
		\end{itemize}
		
	\pagebreak
	\section{Scene graph}
		\subsection{Node}
			
				
		\subsection{Texture}
			The texutre class encapsulates a texture, which resides on the gpu. There are two functions, which allow the creation of either 2D or 3D textures.			
			\begin{lstlisting}
				class Texture
				{
				public:
					explicit Texture();
					~Texture();
					
					/**
					* \brief Create a 2D texture.
					* \return True if creation was successful, false otherwise.
					*/
					bool textureFromFile(const char* path, const std::string& type);
					
					/** 
					* \brief Create a 3D texture.
					* \return True if creation was successful, false otherwise.
					*/
					bool cubeMapFromFile(const char* path, const std::string& type);
				};	
			\end{lstlisting}
		
		\subsection{Material}
			A \textbf{Material} contains a number of properties, which defines the color of an object. These properties include:
			\begin{itemize}
				\item Ambient, Diffuse, Specular colors
				\item Shininess factor for highlights
				\item Textures
			\end{itemize}
			\begin{lstlisting}
				
			\end{lstlisting}	
		
		\subsection{Geometry}
			A geometry object encapsulates various geometry data, including:
			\begin{itemize}
				\item Vertices
				\item Normals
				\item Texture coordinates
			\end{itemize}
			It is responsible for memory managment on the GPU regarding geometry data.
			
			There already are specific geometry implementations for common geometry types:
			\begin{description}
				\item[3D]: 
					\begin{packed_itemize}
						\item CubeGeometry
						\item SphereGeometry
						\item CylinderGeometry
					\end{packed_itemize}
				\item[2D]:
					\begin{packed_itemize}
						\item CircleGeometry
						\item RectangleGeometry
					\end{packed_itemize}
			\end{description}
			If these predefined types aren't sufficicent, you can create your own geometry types. This is done by deriving from \textbf{StaticGeometry}.
			
			
		\subsection{Shader}
			The \textbf{Shader} class can load, compile and link shader programs. The shading language used is GLSL, since OpenGL is being used.
		
		\subsection{Drawable}
			A drawable groups together
			\begin{packed_itemize}
				\item Geometry,
				\item Material and
				\item Shader
			\end{packed_itemize}
		
			Drawables can be added to scene graph \textbf{Node}s so they are rendered.
		
	\pagebreak
	\section{Rendering}
		
		\subsection{Frame}
		
		\subsection{Camera}
		
	
	\pagebreak
	\section{Lighting}
		TBA
		
		
		
	\pagebreak
	\section{Input}
	Input is provided through the GLFW library.
		
		\subsection{Input events}
		Input events are always bound to one window. There are multiple sources, which can generate input events:
		\begin{packed_itemize}
			\item \textbf{Mouse}
			\item \textbf{Keyboard}
			\item \textbf{Window}
			\item \textbf{Joystick}
		\end{packed_itemize}
		An \textbf{InputEvent} consists of an \textbf{InputType} and the \textbf{EventData}. The \textbf{EventData} holds all different kinds of events, of which only one can be active at a time due to it being a union. The \textbf{InputType} describes, which part of the \textbf{EventData} is active. All other elements inside \textbf{EventData} are invalid!
		
		I.e. if the \textbf{InputType} is \textit{key\_pressed}, only the KeyboardEvent inside the \textbf{EventData} is active.

		\subsection{EventSystem}
		An event system is always linked to exactly one window. 
		\begin{description}
			\item[DeviceRegistry:] \hfill\\
				Contains input devices (mouse, keyboard, joystick etc.). Those devices can be activated to generate input events.
			\item[InputObservable:] \hfill\\
				\textbf{InputListener} can be registered at this component, which will then be notified about incoming \textbf{InputEvents}.
			\item[InputListener:] \hfill\\
				Input listener contain a number of callbacks, which can be added by the user. Input listener have to be registered at an \textbf{InputObservable} located in a window, to be informed about input events.	
		\end{description}
		
			
	
	\pagebreak
	\section{GUI}
	The GUI is built using the Agui-library with custom backends for OpenGL.
	
		\subsection{Initialization}
		Before the GUI library-part can be used, it has to be initialized. To do so, use the static function:
		\begin{lstlisting}
			void ec::MiniAgui::init();
		\end{lstlisting}
		
		\subsection{Backend}
		
		\subsection{title}
		
		
	\pagebreak
	
	
	
	
	\pagenumbering{alph}
	\newpage
	\appendix
	\listoffigures
	\listoftables
	
	% TODO: glossary
	
\end{document}